\documentclass[letterpaper, onecolumn, draftclsnofoot, 10pt, compsoc]{IEEEtran}
	\usepackage{graphicx}
	\usepackage[hyphens]{url}
	\usepackage{setspace}
	\usepackage{tabu}
	\usepackage{float}
	\usepackage{pdfpages}
	
	\usepackage{listings}
	\usepackage{xcolor}
	\definecolor{codegreen}{rgb}{0,0.6,0}
	\definecolor{codegray}{rgb}{0.5,0.5,0.5}
	\definecolor{codepurple}{rgb}{0.58,0,0.82}
	\definecolor{backcolour}{rgb}{0.95,0.95,0.92}
	
	\lstdefinestyle{mystyle}{
	  backgroundcolor=\color{backcolour}, commentstyle=\color{codegreen},
	  keywordstyle=\color{magenta},
	  numberstyle=\tiny\color{codegray},
	  stringstyle=\color{codepurple},
	  basicstyle=\footnotesize,
	  breakatwhitespace=false,         
	  breaklines=true,                 
	  captionpos=b,                    
	  keepspaces=true,                 
	  numbers=left,                    
	  numbersep=5pt,                  
	  showspaces=false,                
	  showstringspaces=false,
	  showtabs=false,                  
	  tabsize=2
	}
	
	\lstset{style=mystyle}
	
	\newcommand{\cred}[1]{{\color{red}#1}}
	\newcommand{\cblue}[1]{{\color{blue}#1}}
	
	\usepackage{geometry}
	\geometry{textheight=9.5in, textwidth=7in}
	
	% 1. Fill in these details
	\def \CapstoneTeamName{			The Polycules}
	\def \CapstoneTeamNumber{		63}
	\def \GroupMemberOne{			Joshua Lioy}
	\def \GroupMemberTwo{			Corynna Park}
	\def \GroupMemberThree{			Jackson Wells}
	\def \CapstoneProjectName{		Molecules in 3D?! And in color!? That I can hold in my hand? No way!!!}
	\def \CapstoneSponsorCompany{	Oregon State University College of Science Department of Biochemistry and Biophysics}
	\def \CapstoneSponsorPerson{	Dr. Victor Hsu}
	
	% 2. Uncomment the appropriate line below so that the document type works
	\def \DocType{	%Problem Statement
					%Requirements Document
					%Technology Review
					%Design Document
					Final Report
					}
				
	\newcommand{\NameSigPair}[1]{\par
	\makebox[2.75in][r]{#1} \hfil 	\makebox[3.25in]{\makebox[2.25in]{\hrulefill} \hfill		\makebox[.75in]{\hrulefill}}
	\par\vspace{-12pt} \textit{\tiny\noindent
	\makebox[2.75in]{} \hfil		\makebox[3.25in]{\makebox[2.25in][r]{Signature} \hfill	\makebox[.75in][r]{Date}}}}
	% 3. If the document is not to be signed, uncomment the RENEWcommand below
	\renewcommand{\NameSigPair}[1]{#1}
	
	%%%%%%%%%%%%%%%%%%%%%%%%%%%%%%%%%%%%%%%
	\begin{document}
	\begin{titlepage}
		\pagenumbering{gobble}
		\begin{singlespace}
			\includegraphics[height=4cm]{coe_v_spot1}
			\hfill 
			% 4. If you have a logo, use this includegraphics command to put it on the coversheet.
			%\includegraphics[height=4cm]{CompanyLogo}   
			\par\vspace{.2in}
			\centering
			\scshape{
				\huge CS Capstone \DocType \par
				{\large\today}\par
				\vspace{.5in}
				\textbf{\Huge\CapstoneProjectName}\par
				\vfill
				{\large Prepared for}\par
				\Huge \CapstoneSponsorCompany\par
				\vspace{5pt}
				{\Large\NameSigPair{\CapstoneSponsorPerson}\par}
				{\large Prepared by }\par
				Group\CapstoneTeamNumber\par
				% 5. comment out the line below this one if you do not wish to name your team
				\CapstoneTeamName\par 
				\vspace{5pt}
				{\Large
					\NameSigPair{\GroupMemberOne}\par
					\NameSigPair{\GroupMemberTwo}\par
					\NameSigPair{\GroupMemberThree}\par
				}
				\vspace{20pt}
			}
			\begin{abstract}
				This document is the final report for the Molecules in 3D project. 
				It consists of a project introduction, all the documentation created for the project, as well as the final poster presented at Expo and each group member's conclusions regarding how the project went.
				The purpose of this document is to show what was done, why it was done, and how, in hopes that future audiences may be able to benefit from this project.   
			\end{abstract}     
		\end{singlespace}
	\end{titlepage}
	\newpage
	\pagenumbering{arabic}
	
	\newpage
	\tableofcontents
	% 7. uncomment this (if applicable). Consider adding a page break.
	\newpage
	%\listoffigures
	%\listoftables
	\clearpage
	
	% 8. now you write!https://www.overleaf.com/16812618wkzfzdwqrsmc#
	\begin{singlespace}
	
	\section{Introduction to Project}
	\subsection{Who requested it?}
	This project was requested by Dr. Victor Hsu, a professor in Oregon State University's Department of Biochemistry and Biophysics.
	\subsection{Why was it requested?}
	The project was requested to help enhance 3D printing as a teaching aid through adding a multi-color functionality. 
	Currently, OSU's DBB 3D prints single-colored molecular models to aid students in learning the structure and function of different molecules.
	To show different atom types and surfaces, they then go and hand-paint the single-colored models, which can end up being a tedious and difficult task.
	As a solution, OSU has purchased a multi-filament interface which we were then tasked with integrating it with the 3D printer and using various software to produce multi-colored models printed with up to 4 different filaments.
	\subsection{What is its importance?}
	This project's importance lies in its academic value, but it can also benefit those who wish to 3D print multi-colored objects in general. 
	Those who teach biological and/or chemical subjects will be able to better represent important molecular structures and concepts.
	Those who are tactile and visual learners will be able to learn more effectively through these larger than life, colored models.
	Finally, those who 3D print for non-academic purposes will have a new, fun feature to try out on their own and enhance their hobby.
	\subsection{Who was your client?}
	Our client was Dr. Victor Hsu.
	\subsection{Who are the members of your team?}
	Our team consisted of Joshua Lioy, Corynna Park, and Jackson Wells.
	\subsection{What were their roles?}
	Joshua managed the hardware side of the project, while Jackson took care of development on the software side. Corynna handled any leftover tasks, which mostly involved documentation, and was also the team captain. 
	\subsection{What was the role of the client?}
	The client's role was to provide us with the necessary hardware, specifically the 3D printer, multi-filament interface, computer, lab space, and filaments. 
	He also aided in defining the project's requirements and was a resource for questions when needed.
	
	\section{Requirements Document}
	\subsection{Original Requirements Document}
	\includepdf[pages=-]{capstone_requirements.pdf}
	\subsection{Changes to Requirements Document} % What new requirements were added? What existing requirements were changed? What existing requirements were deleted? Why?
	\begin{itemize}
	\item The unnecessary definitions for CIF, WRL, and OBJ were removed, as they were no longer relevant to the project.
	\item The references to 3MF and CIF were removed so that the file types for the workflow would be restricted to STL, for 3D printable files to be used with the slicing software, and PDB, for biological type files to be used with the conversion software.
	This design change was due to the change in our design, which is addressed in the Changes to Design Document section.
	\item All instances of Blender throughout the documentation were updated to Cura, as Cura was used as the project's slicing software.
	\item All instances of PyMol were removed, as Chimera was used as the project's conversion software.
	\item A section labeled "Object Requirements" was added to explicitly state the primary goal of our project which was to produce multicolored molecular objects, which we accidentally assumed to be implied in the original draft and omitted.
	\end{itemize}
	\subsection{Final Gantt Chart}
	\includepdf[pages=-]{new_gantt.pdf}
	
	\section{Design Document}
	\subsection{Original Design Document}
	\includepdf[pages=-]{capstone_design_document.pdf}
	\subsection{Changes to Design Document} %What design aspects were changed, deleted or added? Why?
	\begin{itemize}
	\item Definitions for STL and PDB were added, as these are the two file types used by our workflow.
	\item A design change that was made was not restricting users to only being able to use 3MF files with the slicing software.
	Such a restriction would force users to perform the extra step of converting a potentially already Cura-compatible file as well as add an additional application or piece of software to the workflow.
	Also, since the 3MF file type is relatively new, this could present challenges in itself, such as lack of interoperability with other software and computing platforms.
	Cura is compatible with STL, OBJ, and 3MF files; the proposal was to change the project's design to accept any of these three file types.
	Ultimately, we ended up using only .STL files in our implementation, as we found that it was relatively simple to create colored 3D data using multiple STL files, and our conversion software could easily convert biological PDB files into STL files.
	\item  The hardware communication section was removed, as it did not accurately reflect the design, but was addressed in the system operation section.
	\item In the conversion software section, all instances of PyMol were removed and replaced with Chimera, which is the conversion software that the group had decided to utilize in the workflow.
	Also PDB was determined to be the file type for biological data.
	\item The system operation section was updated and now has the workflow sending information between the slicing software and the 3D printer through SD card. 
	
	\end{itemize}
	\section{Tech Review}
	\includepdf[pages=-]{capstone_tech_review.pdf}
	
	\section{Weekly Blog Posts}
	\subsection{Joshua's Blog Posts}
	\subsubsection{Fall Term}
\paragraph{Week 01}
\begin{itemize}
\item This week I submitted my choices for potential CAPSTONE projects. 
\end{itemize}
\paragraph{Week 02}
\begin{itemize}
\item Had the initial meeting with group members and Dr. Hsu and covered the basics of the projects. 
\end{itemize}
\paragraph{Week 03}
\begin{itemize}
\item Worked on and submitted the rough draft of the problem statement. 
\item Had a second meeting with Victor to tour the lab and the equipment.  
\end{itemize}
\paragraph{Week 04}
\begin{itemize}
\item Met with group for finalize the problem statement and turned in the final draft. 
\item Started messing with the 3d printer. 
\end{itemize}
\paragraph{Week 05}
\begin{itemize}
\item Met with Daniel on Tuesday and worked on and submitted the rough draft of the requirements document. 
\end{itemize}
\paragraph{Week 06}
\begin{itemize}
\item This week we met with Daniel and completed the final draft of the requirements document. 
\end{itemize}
\paragraph{Week 07}
\begin{itemize}
\item The group met in the lab to go over some of the tech review components. 
\end{itemize}
\paragraph{Week 08}
\begin{itemize}
\item Worked on the rough draft of the tech review. 
\end{itemize}
\paragraph{Week 09}
\begin{itemize}
\item Worked on and turned in final draft of tech review. 
\end{itemize}
\paragraph{Week 10}
\begin{itemize}
\item Had the weekly meeting with Daniel.
\item Worked on and submitted the final draft of the design document. 
\end{itemize}

\subsubsection{Winter Term}
\paragraph{Week 01}
\begin{itemize}
\item Had the initial Capstone meeting for winter term on Tuesday. 
\end{itemize}
\paragraph{Week 02}
\begin{itemize}
\item Had the initial meeting with Daniel for the Winter term. 
\item Got a hold of the Palette and removed the filament blockage. 
\end{itemize}
\paragraph{Week 03}
\begin{itemize}
\item Met with Daniel and finished the basic setup of the Palette. 
\end{itemize}
\paragraph{Week 04}
\begin{itemize}
\item Met with Daniel on Tuesday.
\item Continued working on the Palette.
\item Met with group on Friday.
\end{itemize}
\paragraph{Week 05}
\begin{itemize}
\item Met with Daniel. 
\item Successfully got the first multicolored print. 
\item Started printing with the water molecule files that were sent by Victor. 
\item Unable to finish the print due to support issues. However the transition tower did successfully adhere to the print bed. 
\end{itemize}
\paragraph{Week 06}
\begin{itemize}
\item Worked on and submitted the midterm progress report. 
\end{itemize}
\paragraph{Week 07}
\begin{itemize}
\item Continued working on the project and troubleshooting any remaining adhesion issues with the printer. 
\item Met with Daniel on Tuesday.
\end{itemize}
\paragraph{Week 08}
\begin{itemize}
\item Continued working on the project and troubleshooting any remaining adhesion issues with the printer.
\item Met with Daniel on Tuesday.
\end{itemize}
\paragraph{Week 09}
\begin{itemize}
\item Got the adhesion issues mostly solved and successfully printed the first molecule. 
\item Met with Daniel and discussed the demo for the next week. 
\end{itemize}
\paragraph{Week 10}
\begin{itemize}
\item Unable to meet with Daniel due to a miscommunication. 
\item Sent Daniel a video demo of the software as requested.
\end{itemize}

\subsubsection{Spring Term}
\paragraph{Week 01}
\begin{itemize}
\item Continued working on the project. 
\item Setup meeting time with Daniel.
\end{itemize}
\paragraph{Week 02}
\begin{itemize}
\item Had the initial meeting with Daniel for the Spring term. 
\end{itemize}
\paragraph{Week 03}
\begin{itemize}
\item Met with Daniel. 
\item Worked on four color print.
\end{itemize}
\paragraph{Week 04}
\begin{itemize}
\item Met with Daniel. 
\item Worked on four color print.
\end{itemize}
\paragraph{Week 05}
\begin{itemize}
\item No meeting with Daniel this week.
\item Worked on midterm progress report.
\end{itemize}
\paragraph{Week 06}
\begin{itemize}
\item Made successful print for expo.
\item No meeting with Daniel this week.
\end{itemize}
\paragraph{Week 07}
\begin{itemize}
\item Expo
\end{itemize}
\paragraph{Week 08}
\begin{itemize}
\item No TA meeting this week. 
\item Went to the Wednesday class session. 
\end{itemize}
\paragraph{Week 09}
\begin{itemize}
\item No TA meeting this week.
\item Emailed Victor for final hand-off. 
\end{itemize}
\paragraph{Week 10}
\begin{itemize}
\item Met with Victor to hand-off project on Friday.
\item Met with Daniel on Friday for final TA meeting. 
\end{itemize}
	\subsection{Corynna's Blog Posts}
	\subsubsection{Fall Term}
\paragraph{Week 01}
\begin{itemize}
\item Plans: View and submit bids for Capstone project. 
\item Problems: Need to decide top 5 projects and 2 bottom projects I'd like to work on.
\item Progress: Submitted my preferences Thursday. 
\item Summary: Read project proposals and submitted my preferences. 
\end{itemize}
\paragraph{Week 02}
\begin{itemize}
\item Plans: Get group and project assignment.
\item Problems: Need to contact client and begin problem statement draft for project. 
\item Progress: Got in contact with client, started Problem statement draft.
\item Summary: Met group and got project/client assignment, mailed contact, met client for the first time, and began problem statement draft. 
\end{itemize}
\paragraph{Week 03}
\begin{itemize}
\item Plans: Finish and submit individual problem statement. 
First TA meeting. 
Meet with client to get lab access so we can begin getting familiar with 3D printer. 
Get github repo up and put problem statement on it.
\item Problems: Need to get familiar with the established workflow for 3D printing single-color objects.
Need to set up github repo and begin putting all documentation inside.
\item Progress: Completed draft of Problem Statement and submitted PDF to teach. 
Met with client and learned about current workflow, software being used, 3D printed objects.
\item Summary: Completed individual assignment, had first TA meeting, met with client and got lab access, and set up Github repo.
\end{itemize}
\paragraph{Week 04}
\begin{itemize}
\item Plans: Consolidate individual problem statements into one group problem statement to be sent to client and approved.
Meet with TA.
\item Problems: Group Problem Statement due.
Need computer access in ALS 2157 lab to use 3D printer.
Need access to group OneNote.
\item Progress: Completed Group Problem Statement and emailed it to client for approval.
Client responded  with some revisions; made the revisions and resent for approval.
Met with TA.
Client approved our document.
Emailed Kirsten for access to group OneNote.
\item Summary: Completed and sent Group Problem Statement to client, got client approval, we now have computer and lab access and can begin working with the 3D printer.
\end{itemize}
\paragraph{Week 05}
\begin{itemize}
\item Plans: Begin work on Project Requirements document (rough draft due this Friday on OneNote).
\item Problems:  Need to define project requirements, aka what criteria our project be graded on.
\item Progress: Met with group to begin working on document.
\item Summary:  Had TA meeting, began drafting problem requirements document.
\end{itemize}
\paragraph{Week 06}
\begin{itemize}
\item Plans: Project Requirements document.
\item Problems: Need to establish project requirements and send to client for approval.
\item Progress: Completed all our goals for the week.
\item Summary: Finished project Requirements Document and sent to client. 
Client responded with some corrections. 
Made necessary corrections and resubmitted to client.
\end{itemize}
\paragraph{Week 07}
\begin{itemize}
\item Plans: Begin work on tech review.
\item Problems: Need to define 9 different technologies for our project.
\item Progress: Met in the lab to go over established 3D printing workflow and discuss project roles.
\item Summary: Divided up 9 technologies, 3 per person, to begin work on individual tech reviews. 
Rough draft due in class on Tuesday, 11/14.
\end{itemize}
\paragraph{Week 08}
\begin{itemize}
\item Plans: Complete rough draft of individual tech review and bring to class for peer review, TA meeting.
\item Problems: Need to define 3 options for the three pieces of technology assigned to me (filament type, file type, slicing software) and determine which option is best for the project.
\item Progress: Began research on three pieces of technology and possible options.
\item For my portion of the tech review, the pieces and options will be as follows:
	\begin{itemize}
    \item Filament type: PLA, ABS, or Nylon
	\item File type: STL, OBJ, or 3MF
	\item Slicing software: Voxelizer, Simplify3D, or Cura
    \end{itemize}
\end{itemize}
\paragraph{Week 09}
\begin{itemize}
\item Plans: No class Tuesday, finalize draft of tech review.
\item Problems: Need to edit and finish my individual tech review and send digital PDF to Kirsten by 11:59 PM Tuesday.
\item Progress: Finished research and document and submitted. 
\item Summary: Finished and submitted individual technology review. 
Next, design document.
\end{itemize}
\paragraph{Week 10}
\begin{itemize}
\item Plans: Complete and submit Design Doc Friday, 11:59PM.
\item Problems: Need to be able to prioritize time and complete all tasks for all classes.
\item Progress: Worked on design doc both with team and individually. 
Was able to divide up tasks to make the assignment more manageable for us individually.
\item Summary: Completed and submitted design doc by due date, survived dead week.
\end{itemize}
\subsubsection{Winter Term}
\paragraph{Week 01}
\begin{itemize}
\item Plans: Begin the term.
\item Problems: 8AM meeting time T and R, however, only meeting a few times throughout the term.
\item Progress: Not much accomplished this week; first week back. Adjusting to the new term.
\item Summary: Began Winter term. 
Not much happened in terms of project work this week. 
Will be meeting with the group next week as well as have the first TA meeting of the term.
\end{itemize}
\paragraph{Week 02}
\begin{itemize}
\item Plans: Get in contact with client to get Palette+ and feedback/verification on documents. 
Jackson to start python pipeline portion of the project.
\item Problems: No client response since beginning of Christmas break.
\item Progress: Client responded on Tuesday about getting the Palette+ to the group as well as looking at the design document.
\item Summary: We are limited on the progress we can make until we get the Palette+ installed in the lab. 
We have client correspondence and will follow up to get the Palette+ and Design Doc verification. 
Beginning next TA meeting, the TA would like to see some physical representation of project progress every TA meeting.
\end{itemize}
\paragraph{Week 03}
\begin{itemize}
\item Plans:  TA meeting, begin task of looking into how to convert/utilize 3MF files.
\item Problems: Still unable to really begin working with the hardware and interfacing the software with the hardware due to issues getting the Palette+ and Zmorph to produce a successful print.
\item Progress: Began research Thursday.
\item Summary: Will need to keep looking into how we will be using 3mf file type. 
Hopefully the issue with the hardware will be resolved soon.
\end{itemize}
\paragraph{Week 04}
\begin{itemize}
\item Plans:  Meet with group and play catch up after being sick.
\item Problems: I was sick.
Still having hardware (adhesion between the bed and the filament).
\item Progress: Group meeting to go over developments, progress with the hardware, tasks.
\item Summary: Still limited on what we are able to do because of hardware issues. 
\end{itemize}
\paragraph{Week 05}
\begin{itemize}
\item Plans: TA meeting, not much else for this week. 
Begin poster and progress report.
\item Problems: Still having adhesion issues with 3D printer. 
\item Progress: Met with TA, looked at assignment requirements.
\item Summary: Feeling a little at a standstill since we are still having hardware issues. 
But we're getting closer. 
Will need a UML diagram I should look into that.
\end{itemize}
\paragraph{Week 06}
\begin{itemize}
\item Plans:  Submit preliminary poster and midterm progress report.
\item Problems: N/A
\item Progress: Finished poster and emailed to TA. 
Submitted progress report.
\item Summary: Submitted assignments. 
Will need to have a deliverable by week 8 TA meeting.
\end{itemize}
\paragraph{Week 07}
\begin{itemize}
\item Plans: Update design and requirements doc, work on alpha build for next TA meeting.
\item Problems: N/A
\item Progress: Updated docs on Friday.
\item Summary: Getting close to having all requirements done. 
We need to make sure to have a deliverable by next TA meeting.
\end{itemize}
\paragraph{Week 08}
\begin{itemize}
\item Plans: TA Meeting.
\item Problems: N/A
\item Progress: Presented TA with alpha demo, approx. 7/9 requirements met.
\item Summary: Need to update the design and requirements documents, which will bring us closer to having fulfilled all our requirements.
\end{itemize}
\paragraph{Week 09}
\begin{itemize}
\item Plans: TA meeting.
\item Problems: N/A 
\item Progress: Met with TA and scheduled next meeting to be in the lab to demo.
\item Summary: Need to get set up for our beta demo next TA meeting. 
Sent TA updated design and requirements docs.
\end{itemize}
\paragraph{Week 10}
\begin{itemize}
\item Plans: Demo for the TA in the lab in place of TA meeting.
\item Problems: TA did not make it to meeting.
\item Progress: Have implemented all requirements, now just need to debug.
\item Summary: Unable to do lab demo for the TA but all project requirements have been met. 
Will send a video demo instead.
Winter progress report due next Wednesday.
\end{itemize}
\subsubsection{Spring Term}
\paragraph{Week 01}
\begin{itemize}
\item Plans: Schedule TA meeting.
Register for expo.
\item Problems: N/A
\item Progress: Registered for Expo on Wednesday, April 4.
\item Summary: Submitted timeslot for TA meeting time on Friday from 3-5. 
Registered for Expo.
\end{itemize}
\paragraph{Week 02}
\begin{itemize}
\item Plans: TA Meeting on Friday.
Not much left to do besides expo poster and release forms (May 1st).
\item Progress: Had TA meeting and discussed limitations of hardware and future expectations.
\item Problems: N/A
\item Summary: Hardware still limiting us on ability to consistently produce successful prints. 
Need to finalize poster and submit, and turn in release forms.
\end{itemize}
\paragraph{Week 03}
\begin{itemize}
\item Plans: TA Meeting on Friday.
\item Progress: N/A
\item Problems: N/A
\item Summary: Hardware still being serviced, still being a problem. 
TA wants one successful print and then we will be done.
\end{itemize}
\paragraph{Week 04}
\begin{itemize}
\item Plans: TA Meeting on Friday.
Finalize poster and submit through canvas assignment, and submit for printing. 
Turn in release forms.
\item Progress: Completed all tasks listed in the plan.
\item Problems: N/A
\item Summary: Submitted final draft of poster through canvas on Friday. 
Emailed signed release forms on Friday.
TA meeting: Have code ready by Friday, May 11th. 
Do not need checkoff as we had that at the end of last term. 
Further TA meetings will be optional; email for appointment.
\end{itemize}
\paragraph{Week 05}
\begin{itemize}
\item Plans: Work on Spring term Progress Report and Presentation due Sunday, May 6th.
\item Progress: Began working on report and presentation.
\item Problems: N/A
\item Summary: Submitted progress report by the deadline. 
No TA meeting this week.
\end{itemize}
\paragraph{Week 06}
\begin{itemize}
\item Plans: Class on Wednesday.
\item Progress: Attended class.
\item Problems: N/A
\item Summary: Had class on Wednesday and were given reminders for Expo (next Friday).
\end{itemize}
\paragraph{Week 07}
\begin{itemize}
\item Plans: Expo on Friday!
\item Progress: Prepared for Expo.
\item Problems: N/A
\item Summary: Presented our project at Expo.
\end{itemize}
\paragraph{Week 08}
\begin{itemize}
\item Plans: Go to class Wednesday.
\item Progress: Attended class.
\item Problems: N/A
\item Summary: Class on Wednesday, reminded to do final project handoff and about various upcoming assignments.
\end{itemize}
\paragraph{Week 09}
\begin{itemize}
\item Plans: Contact client to schedule final handoff (in form of a demo).
Upcoming assignments: Peer review, final presentation (Due Friday, June 8), final report (Due Tuesday, finals week).
\item Progress: Emailed client to request meeting on Wednesday, May 30th.
\item Problems: No client response by the end of the week.
\item Summary: Emailed client. 
No response yet by the end of the week, will email again if nothing heard over the weekend. 
Completed peer review this week.
\end{itemize}
\paragraph{Week 10}
\begin{itemize}
\item Plans: Email client again to request demo.
Begin work on final presentation (due Friday) and final report (due next Tuesday).
Class Wednesday.
\item Progress: Client responded Thursday with his availability. 
Scheduled demo for Friday at 2:30PM.
Attended class Wednesday, which consisted of students giving feedback regarding their time here at OSU.
\item Problems: N/A
\item Summary: Presented final project to client in the form of an in-lab demo during which the software and hardware components were presented.
Began, completed, and submitted our final presentation.
Last assignment, final report, due Tuesday.
\end{itemize}
	\subsection{Jackson's Blog Posts}
	\subsubsection{Fall Term}
\paragraph{Week 01}
\begin{itemize}
\item Selected preferences from project list
\item Submitted resume and personal biography.
\end{itemize}
\paragraph{Week 02}
\begin{itemize}
\item Met with client.
\item Met with group.
\item Having some issues with the scheduling of our weekly TA meeting, all group members are available at different times and I have class during every time provided by the TA. I emailed the TA in hopes of working through the issue, which we currently are. 
\end{itemize}
\paragraph{Week 03}
\begin{itemize}
\item Worked individually on the problem statement.
\item Met with the client again, this time in the lab where the 3D printers are housed. 
\item We established a GitHub page and invited the client and another individual 
In the department that has a stake in the project.
\end{itemize}
\paragraph{Week 04}
\begin{itemize}
\item Met with TA
\item Consolidated problem statement.
\item The client had 3-4 minor suggestions for edits, which were corrected and submitted.
\end{itemize}
\paragraph{Week 05}
\begin{itemize}
\item Requirements document partly drafted.
\item Some prints completed using the printer.
\item Set future times to continue work on the requirements document.
\end{itemize}
\paragraph{Week 06}
\begin{itemize}
\item Requirements document finalized.
\item Gantt chart created for addition to the requirements document.
\item The client approved our document, with a few grammatical corrections.
\end{itemize}
\paragraph{Week 07}
\begin{itemize}
\item Technology Review partially written.
\item Our group met to get hands on time with the printer, and to get a set list of technological components that are included in our project. We divided up the components equally and proceeded from there.
\end{itemize}
\paragraph{Week 08}
\begin{itemize}
\item Gained insight from peer review of technology review.
\item Continued work on the technology review.
\end{itemize}
\paragraph{Week 09}
\begin{itemize}
\item Finalized and submitted the technology review.
\end{itemize}
\paragraph{Week 10}
\begin{itemize}
\item Began and finalized the design document with the group.
\end{itemize}

\subsubsection{Winter Term}
\paragraph{Week 01}
\begin{itemize}
\item Communication was reestablished with the TA and a meeting was set for the following week.
\end{itemize}
\paragraph{Week 02}
\begin{itemize}
\item Met with the TA to provide feedback on the previous term and overview the current term.
\item I was given the goal of implementing the input handling functions by next week's meeting.
\end{itemize}
\paragraph{Week 03}
\begin{itemize}
\item Input handling was implemented.
\item Skeleton code was produced to outline the remaining aspects of the work-flow's software.
\end{itemize}
\paragraph{Week 04}
\begin{itemize}
\item Met with the group and the TA on separate occasions. Group meeting served to summarize the TA meeting due to sickness related absences.
\item A new goal of investigating the process for handling of PDB files was set for completion by the next TA meeting.
\item The hardware aspect of the work-flow was presenting numerous issues, preventing print production.
\end{itemize}
\paragraph{Week 05}
\begin{itemize}
\item I was able to replicate the O and H STL files provided by the client! This is a big step forward in terms of project development. We should now in theory be able to produce polychromatic prints of biological molecules.
\item The software tools to be utilized in the work-flow were locked in, and are now ready for implementation.
\item Results were shared at the TA meeting. Conveniently timed, we have also been able to produce a polychromatic test print using the new hardware.
\end{itemize}
\paragraph{Week 06}
\begin{itemize}
\item System commands were investigated and implemented (for Chimera and Cura) to execute the dependent printing software tools from the work-flow script.
\item The next task of input file validation was assigned.
\end{itemize}
\paragraph{Week 07}
\begin{itemize}
\item Met with the TA.
\item Development/investigation of file validation continued.
\end{itemize}
\paragraph{Week 08}
\begin{itemize}
\item Added Chroma to the pipeline (system calls) and finalized file validation in time for the alpha deliverable.
\item Met with Daniel on Tuesday to show our alpha deliverable. We went over the requirements and design documents and discussed revisions. 
\end{itemize}
\paragraph{Week 09}
\begin{itemize}
\item Met with the TA and set a time to exhibit the beta demo of the work-flow.
\item Implemented further revisions to the design and requirements documents
\end{itemize}
\paragraph{Week 10}
\begin{itemize}
\item A miscommunication occurred between the group and the TA and the meeting was not conducted. The TA believed we would first meet in Kelley and then walk to the lab. The group met in the lab.
\item The TA requested a video be recorded to exhibit beta performance and one was supplied.
\end{itemize}

\subsubsection{Spring Term}
\paragraph{Week 01}
\begin{itemize}
\item A Friday meeting time was established with the group and TA.
\item Having provided a beta demo last term a week 1 meet was deemed unnecessary.
\end{itemize}
\paragraph{Week 02}
\begin{itemize}
\item Met with the TA.The meeting was brief and the group had a discussion following the TA meeting to get some things sorted out with upcoming assignments.
\item We plan to begin revisions on the poster this weekend and into next week.
\end{itemize}
\paragraph{Week 03}
\begin{itemize}
\item Met with the TA. General check in on progress.
\end{itemize}
\paragraph{Week 04}
\begin{itemize}
\item Met with the TA. General check in on progress.
\end{itemize}
\paragraph{Week 05}
\begin{itemize}
\item Began and finalized work on the midterm progress report.
\end{itemize}
\paragraph{Week 06}
\begin{itemize}
\item Added support for Windows and OSX operating systems too the software work-flow.
\item Finalized file validation and added comments throughout the code.
\item Successfully produced a polychromatic print of a molecule.
\end{itemize}
\paragraph{Week 07}
\begin{itemize}
\item Presented our project at Expo
\end{itemize}
\paragraph{Week 08}
\begin{itemize}
\item Cool down week.
\end{itemize}
\paragraph{Week 09}
\begin{itemize}
\item Contacted the client in regards to the project hand off, no response.
\end{itemize}
\paragraph{Week 10}
\begin{itemize}
\item Conducted the product hand off with the client. Provided overviews of the hardware and software aspects of the project.
\item Met with the TA. Discussed the hand off and remaining work to be completed before the term ended.
\end{itemize}
	
	\section{Final Poster}
	\includepdf[pages=-, angle=90]{team_63_expo_poster.pdf}
	
	\section{Project Documentation}
	\subsection{Project Structure} %See the section 6 of the assignment description
	The delivered python script was structured in a user-friendly fashion. Self-explanatory variable names were utilized and comments were included throughout the code, in the case that the user wishes to delve into the code and potentially make changes. The script itself makes use of an additional configuration file which simply serves to hold the installation locations of dependent software tools. Both the project's script and required configuration file are stored inside the project repository's dev/ directory. \par
	When executed, the script begins by reading the configuration file and proceeds to check the flags supplied in the executing command. The script will exit if an input file is not supplied or if a host operating system is not specified. The supplied input file is then validated, with the script exiting in error if the file fails any of the validation tests. Following validation, the dependent software packages are executed. If a PDB file is input, Chimera will be opened first and will display the supplied file. If an STL file is supplied, Cura will be launched and the supplied file will be displayed. After closing the first executed software tool, the second software tool will be opened and the user will need to manually open the files created in the previous dependent software.
	\subsection{Installation and Running of Software}
	There are four total dependencies that are required for operation of the delivered work-flow script: Python, Chimera, Cura and Chroma. The client required that all software utilized by our script must be open-source and available for download to the general public, so the project was easily portable. We met this requirement and all the of dependent software packages/tools can be downloaded from their respective host organizations.
	\subsubsection{Installation}
	The installation process will differ depending upon the operating system of those wishing to execute our pipeline script, but all software tools are presented with OSX and Windows based installation files. The actual installation process for each dependent will be the same as any other software that was previously installed on the users system. Each package comes complete with a step-by-step GUI to walk the user through the installation progress. Be sure to record the installation paths of Cura, Chimera and Chroma as this information will be required for execution of the project's script. For Windows users only, be sure to check the box specifying the addition of Python to the PATH variable during installation, or else you will have to supply Python's full installation path every time the project's script is executed. In addition to the Python script there are two profile files for Cura and Chroma that will need to be loaded if the software is being used on a different computer.
	\subsubsection{Running The Software}
	To execute the molecule printing pipeline, first obtain all the files contained in the project's GitHub repository (if you're reading this you might already have completed this step). Following acquisition of the repository's files, navigate to the dev\/ directory, open the paths\.conf file using the text editor of your choice and enter the installation location of each of the dependent software tools installed in the previous section. Mac users may not have to make any changes to these lines as the file is by default setup to work with OSX based operating systems (if your installation directory varies from the defaults be sure to include \.app at the end of the application's path). Windows users should be sure to place the path of the \.exe file of each software tool. Save these changes (if any) to the paths\.conf file and exit. You are now ready to execute the project's script. A typical run of the program will have syntax similar to the following: 
	\begin{lstlisting}[language=bash]
	python polyPipeline.py -m -i ../test_files/water.pdb
	\end{lstlisting}
	The script also comes complete with a help menu/usage statement that will be printed when the program is executed using either of the following examples: 
	\begin{lstlisting}[language=bash]
	python polyPipeline.py -h
	\end{lstlisting}
	or
	\begin{lstlisting}[language=bash]
	python polyPipeline.py --help
	\end{lstlisting}
	
	\subsection{Special Requirements}
	For our project the only special requirements we had were the use of the Palette+, the ZMorph, Chimera, Cura, Chroma, and the Python runtime environment.
	
	\section{Recommended Technical Resources} %See section 7 of the assignment description
	\subsection{Helpful Websites}
	\begin{itemize}
	\item To learn more about Ultimaker's Cura and download it from the web: \url{https://ultimaker.com/en/products/ultimaker-cura-software}
	\item For documentation on Mosaic Manufacturing's Palette+ and the Chroma software: \url{https://www.mosaicmanufacturing.com/pages/technology}
	\item Download and learn about UCSF's Chimera: \url{https://www.cgl.ucsf.edu/chimera/}
	\end{itemize}
	\section{Conclusions and Reflections}
	\subsection{Joshua}
	\subsubsection{What technical information did you learn?}
For my portion of the project which was primarily focused on the hardware; thus, the majority of the technical information that I learned was centered around the hardware. I gained a much greater understanding of the intricacies of 3D printing hardware and some of the software tools that interface with said hardware. This included things such as common issues with 3D printers and the relatively early state of multifilament printing at the consumer level.
\subsubsection{What non-technical information did you learn?}
Relative to the amount of technical information that I learned working on this project, I do not think I learned much new non-technical information. For me the structure of the course and nature of our project was not overly conducive to introducing new non-technical information. The one exception to this was preparing and participating in the expo. I had not previously participated in an event of that scale and it was interesting to experience it with my group members.
\subsubsection{What have you learned about project work?}
One point that was reinforced over the course of the project was to ensure that communication with the client was kept open during crucial times during the project such as the initial documentation and the final hand-off of the deliverables at the end of this term.
\subsubsection{What have you learned about project management?}
I learned that it is important to have clearly defined goals at the start of the project and an action plan in place to accomplish those goals.
\subsubsection{What have you learned about working in teams?}
The most important point that was reinforce over the course of the project was maintaining consistent communication with my teammates to ensure that problems are addressed in a timely manner and that collaborative efforts can be made more effective.
\subsubsection{If you could do it all over, what would you do differently?}
If I could do one thing differently it would be to spend more time testing four colored prints. This was one of the things that I wanted to do but was unable to finish before expo.
	\subsection{Corynna}
	\subsubsection{What technical information did you learn?}
Working on this project introduced me more to the concept of 3D printing and the functions of a 3D printer and its parts, as well as the different technologies involved.
I also learned how to research and design a technical project, as well as how to write documentation for one.
\subsubsection{What non-technical information did you learn?}
The non-technical information that I learned was how to adapt a pitch to the possible different technical levels of your audience.
From the in-class exercise to presenting at Expo, being able to analyze your audience and present accordingly is a valuable life skill.
\subsubsection{What have you learned about project work?}
From working on this project I learned that it is good to have lots of documentation so that you know exactly what goals you are trying to accomplish and how.
It is also good to have defined roles for each member.
Throughout the development of the project, these are subject to change, but having a good plan will limit the amount of avoidable issues that often come with working on a project.
\subsubsection{What have you learned about project management?}
Project management requires good prioritizing and planning skills, as well as flexibility.
Throughout the course of this project, what helped was setting goals and having a plan of when to accomplish them by.
Of course, due to unforeseen challenges our timeline varied and we had to remain flexible regarding when we wanted to have certain tasks completed by.
\subsubsection{What have you learned about working in teams?}
Working in teams requires a balance of control and dependency. 
While there are tasks that each member may be best-suited for, a successful team needs a lot of productive communication and a good sense of direction and control.
As a part of Group 63, I found that all were self-sufficient, well-working individuals.
The challenge that comes with that is the lack of communication because each person is able to complete tasks well on his or her own.
And while I do not think this caused a problem for us, for future larger projects it could lead to inconsistencies.
\subsubsection{If you could do it all over, what would you do differently?}
If I could do this project all over again, I would take more initiative during the development phases so that I could be more hands-on and involved, specifically with trouble-shooting the hardware.
I felt that the problems with hardware is what restricted our timeline for the most part, and while we did get to where we needed to be by the deadline, I just wish I had been able to help get past sticking points quicker.
 
	\subsection{Jackson}
	\subsubsection{What technical information did you learn?}
Through the course of this project, I gained technical insights in both the hardware and software aspects of development. In terms of hardware, I obtained a relatively solid background in the process of 3D printing, starting with a structural file to production of the print. I learned the nomenclature associated with the hardware devices, and the intended functionality of particular components, like the extruder or multi-filament interface. In terms of software, I learned new concepts like operating system interoperability and execution of dependent software tools that were not natively accessible from a command-line environment.
\subsubsection{What non-technical information did you learn?}
While the project involved mostly technical subjects, I managed to learn a few skills that can be put to use later in my career. I learned proper communication standards that are to be expected when interacting with a client, and how bargaining might take place in terms of work to be completed on a proposed project. I also gleaned experience in time management and information retrieval for new concepts encountered during development.
\subsubsection{What have you learned about project work?}
In development of the project I discovered the importance of time management. The weekly TA meetings served as a constant reminder that there was work to be completed. At some points I found myself scrambling to meet deadlines which created stress that could have been circumvented through the implementation of better personal scheduling.
\subsubsection{What have you learned about project management?}
I found that documentation plays a massive roll in the management of a project. There were a couple of instances during the three terms where I felt I could have taken more in-depth notes or better recorded completed work for future reference. Including such notes allow for a reduction in the amount of work that must initially be completed when returning to development after a break between terms or in switching from other work from courses being completed concurrently with the capstone series.
\subsubsection{What have you learned about working in teams?}
Our group was efficient, communicated well while also being self-sufficient. Due to the lack of any real problems with the dynamic between members I don't believe there was too much to learn. The weekly TA meetings gave us a way to track each others progress on our respective parts of the project. We maintained a group text message throughout all three terms which provided us a way to quickly get in touch with everyone when issues were encountered. 
\subsubsection{If you could do it all over, what would you do differently?}
If we were to begin the project again from Fall term, I would attempt to bargain more with the client to procure better printing hardware. I believe adhesion and filament issues are what were the most detrimental to development, and that most of these issues could have been avoided if we had access to hardware that was closer to industry standards. With the elimination of these complications we could have considered the outlined stretch goals like the production of a publication or the implementation of a GUI to reduce the need for prior experience in a command-line environment.
	
	\subsection{Appendix 1: Essential Code Listings}
	\subsubsection{Overview}
	Here is the code for our python script as it stood at the end of our project's development. Note that this code is feature complete.
	
	\begin{lstlisting}[language=Python]
	#!/usr/bin/env python
	# coding=UTF-8
	import getopt, sys
	import os
	import subprocess
	import re
	
	inputFile = ""
	chimDir = ""
	curaDir = ""
	chroDir = ""
	
	def launchSliceWindows(inputString):       	#Method for launching Windows slicing software
			cmd = '%s %s' % (inputString,curaDir)
			subprocess.call(cmd,shell=True)
			return
	
	def launchConversionWindows(inputFile):       	#Method for launching Windows conversion software
			cmd = "%s %s" % (inputFile,chimDir)
			subprocess.call(cmd,shell=True)
			return
	
	def launchPrintWindows():                   	#Method for launching Windows printing software
			cmd = "%s" % chroDir         		#create cmd string with wait flag to hold execution
			subprocess.call(cmd,shell=True)
			return
	
	def launchSliceOSX(inputString):		#Method for launching OSX slicing software
		cmd = 'open -W %s -a %s' % (inputString,curaDir)
		subprocess.call(cmd,shell=True)
		return
	
	def launchConversionOSX(inputFile):		#Method for launching OSX conversion software
		cmd = "open -W %s -a %s" % (inputFile,chimDir)
		subprocess.call(cmd,shell=True)
		return
	
	def launchPrintOSX():				#Method for launching OSX printing software
		cmd = "open -W -a %s" % chroDir		#create cmd string with wait flag to hold execution
		subprocess.call(cmd,shell=True)
		return
		
	def validSTL(inputFile):			#method for validation of STL files
		if not os.path.isfile(inputFile):	#checks readability
			print ("File Does Not Exist!\n")	
			sys.exit(1)
		return 1 #file type
	
	
	def validPDB(inputFile):			#method for PDB file validation
		if not os.path.isfile(inputFile):	#checks readability
					print ("File Does Not Exist!\n")
					sys.exit(1)
		endFlag = 0				#flag for detection of required file syntax
		regexp = re.compile(r'[A-Za-z0-9.-]')
		f = open(inputFile, "r")
		for line in f:				#loop over the file
			if "END" in line:		#if END is found
				endFlag = 1
			elif not (regexp.search(line)):	#if invalid characters are found
				print("Invalid File Format!\n")
				sys.exit(2)
		if endFlag:
			return 2
		else:					#if no END was found
			print("Invalid File Format!\n")
			sys.exit(2)
	
	
	def fileValidate(inputFile):			#method that manages validation
		tempType = 0
		ext = inputFile[-4:]		#get extension from file name
		if(ext[0] == "."):
			ext = ext[-3:]
		if(ext == "pdb" or ext == "PDB"):
			tempType = validPDB(inputFile)	#validate the PDB file
		elif(ext == "stl" or ext == "STL"):
			tempType = validSTL(inputFile)  #validate the STL file
		else:					#if file is not PDB or STL
			print ("Invalid File Type!\n")
			sys.exit(1)
		return tempType
	
	
	def usage():		#method for outputing usage information
		print ("\nPolyChromatic 3D Printing Pipeline\n\nDeveloped by:\n\tJoshua Lioy, Corynna Park, Jackson Wells\n\nUsage:\n\tpython polyPipeline.py [Optional Flags] [-w or -m]  -i [input file]\n\nFlags:\n\t-h\tdisplayes usage\n\t-i\tto input file\n\t-v\tfor verbose program output\n\t-w\tsets operating system to Windows\n\t-m\tsets operating system to OSX\n\n")
	
	
	def readConfiguration():				#method for parsing the configuration file
		global chimDir,chroDir,curaDir		
		f = open("paths.conf","r") 
		for line in f:				#loop over config file
			if "Chimera:" in line:	
				chimDir = line.rstrip()
				chimDir = chimDir[8:]
			elif "Cura:" in line:			#store location information
				curaDir = line.rstrip()
				curaDir = curaDir[5:]	
			elif "Chroma:" in line:	
				chroDir = line.rstrip()
				chroDir = chroDir[7:]		
		f.close()
		return
	
	
	def main():
		readConfiguration() 		#parses and stores information from paths.conf
		try:				#sematically identical to getopt in C
			opts, args = getopt.getopt(sys.argv[1:], "hi:vwm", ["help", "input="]) #possible input flags
		except getopt.GetoptError as err:
			print (err)
			usage()
			sys.exit(2)
		verbose = False
		fileInput = False
		OS = -1
		for o, a in opts:
			if o == "-v":
				verbose = True			#prints execution information
			elif o in ("-h", "--help"):
				usage()				#prints usage statement
				sys.exit()	
			elif o in ("-i", "--input"):
				inputFile = a			#supplies the input file
				fileInput = True
			elif o == "-w":				#flag for windows users
				OS = 1
			elif o == "-m":				#flag for Mac users
				OS = 2
			else:					#unknown flags handled here
				assert False, "unhandled option"
		if not fileInput:				#checking that input was made
			print ("\nNo input file specified!\n\nUse -h to display usage statement\n")	
			sys.exit(2)
		else:	#input was supplied
			if(verbose):
						print ("Validating input file\n")
			fileType = fileValidate(inputFile)		#validate input file
		if(verbose):
					print ("File valid!\n")
		if(fileType == 1):			#if the input file is in STL format
			if(verbose):
				print ("Launching CURA\n")
			if(OS == 1):
				launchSliceWindows(inputFile)		
			elif(OS == 2):
				launchSliceOSX(inputFile)
			else:
				print ("Failure to supply an operating system type!\n")
				sys.exit(2)
			if(verbose):
				print ("Launching Chroma\n")
			if(OS == 1):
							launchPrintWindows()
					elif(OS == 2):
							launchPrintOSX()
					else:
							print ("Failure to supply an operating system type!\n")
							sys.exit(2)
			sys.exit(0)	
		elif(fileType == 2):			#if the input file is in PDB format
			if(verbose):
				print ("Launching Chimera\n")
			if(OS == 1):
							launchConversionWindows(inputFile)
					elif(OS == 2):
							launchConversionOSX(inputFile)
					else:
							print ("Failure to supply an operating system type!\n")
							sys.exit(2)
			if(verbose):
				print ("Launching CURA\n")
			if(OS == 1):
							launchSliceWindows("")
					elif(OS == 2):
							launchSliceOSX("")
					else:
							print ("Failure to supply an operating system type!\n")
							sys.exit(2)
			if(verbose):
				print ("Launching Chroma\n")
			if(OS == 1):
							launchPrintWindows()
					elif(OS == 2):
							launchPrintOSX()
					else:
							print ("Failure to supply an operating system type!\n")
							sys.exit(2)
			sys.exit(0)
		else:					#handling cases where fileType value is unknown
			sys.exit(0);
	
	if __name__ == "__main__":
		main()
	\end{lstlisting}
	
	\subsection{Appendix 2: Images}
	\begin{figure}[H]
	\centering
	\includegraphics[scale=0.2]{full_setup}
	\caption{This is the full setup, and shows the rack that holds the filament along with the Palette+ and the printer.}
	\end{figure}
	
	\end{singlespace}
	\end{document}
	